% =====================================================
% ===================== ESPAÑOL =======================
% =====================================================

\cvheader{Juan Diego Figari}{Buenos Aires, Argentina}{juancito.diego@gmail.com}{https://github.com/Yarthax23}

\section*{Perfil}
Estudiante de Ciencias de la Computación (UBA) con experiencia en producto: encuentro lo que falla
usando el software, lo reproduzco y lo arreglo o lo reporto bien documentado.
Trilingüe (español, inglés, francés), con experiencia en atención al cliente y en docencia.

\section*{Experiencia Laboral}
\textbf{Urbanly LLC~--- Ingeniero Junior} \hfill Ene 2026--presente \\
Software de simulación para planificación urbana (CityCompass: TypeScript, React, MapLibre).
\begin{itemize}
  \item \textbf{Debugger de la simulación.} Para un video del producto, el debugger pausaba cada mes en las
        subastas de terrenos y de viviendas, pero no permitía seguir a nadie en particular. Propuse y construí el
        seguimiento: elegir una o varias parcelas o desarrolladores y seguirlos desde la subasta del terreno hasta
        la venta de las viviendas, con un modo que saltea los meses en los que no pasa nada.
  \item \textbf{Red de transporte.} Mientras preparaba un video demo (modificar la red y ver cómo cambia la
        accesibilidad), encontré por qué una línea nueva no tenía efecto: se creaba inactiva por su rango de años.
        También encontré dos errores que dejaban vacía la accesibilidad en colectivo. Los reproduje y los arreglé,
        y agregué botones para activar y desactivar todas las líneas a la vez.
  \item \textbf{Editor de paletas de colores.} Construí el editor de colores y rangos de las capas del mapa,
        iterando con diseño a través de varias revisiones. Después encontré y arreglé un caso en que escribir un
        valor borraba clases sin aviso: una paleta de 9 clases quedaba en 2.
  \item \textbf{Viviendas accesorias (ADUs).} Colaboración con investigadores de Cal Poly Pomona sobre la
        construcción de ADUs en Los Ángeles: el modelo, la validación contra permisos reales y el informe para
        los colaboradores.
\end{itemize}

\noindent \textbf{McDonald's~--- Arcos Dorados S.A.} \hfill 2023--2024
\begin{itemize}
  \item Atención al cliente, manejo de caja y cierre; reconocido como presentador destacado.
\end{itemize}

\noindent \textbf{Profesor de Piano (freelance)} \hfill 2021--2025
\begin{itemize}
  \item Clases particulares adaptadas a cada alumno; gestión autónoma de alumnos y horarios.
\end{itemize}

\section*{Proyectos}
\noindent \textbf{Visualización Geoespacial~--- Subte de Buenos Aires} (\href{https://github.com/Yarthax23/urbanly-subte-map}{GitHub}) \hfill \textit{2026}
\begin{itemize}
  \item Mapa web con MapLibre GL JS y datos abiertos del Gobierno de la Ciudad de Buenos Aires, con un
        pipeline offline que normaliza fuentes en CSV y GeoJSON.
\end{itemize}

\noindent \textbf{Servidor de Chat Event-Driven en C} (\href{https://github.com/Yarthax23/unix-chat-server}{GitHub}) \hfill \textit{2025}
\begin{itemize}
  \item Servidor multi-cliente con sockets UNIX y \texttt{select()}; protocolo en capas y diseño pensado para
        poder depurarlo.
\end{itemize}

\noindent \textbf{Otros:} API REST en Go con PostgreSQL
(\href{https://github.com/Yarthax23/go-todo-api}{repo}); ejercicios de programación funcional y lógica en
Haskell y Prolog (\href{https://github.com/juop12/Exactas-PLP-cNil}{repo}).

\section*{Habilidades}
\begin{itemize}
  \item \textbf{Calidad de producto:} testing exploratorio, reproducción y reporte de bugs, herramientas de desarrollo del navegador
  \item \textbf{Frontend:} TypeScript, React, MapLibre GL, JavaScript, HTML/CSS
  \item \textbf{Colaboración:} Git y GitHub (issues, pull requests, revisión de código), trabajo con equipos de diseño e investigadores externos
  \item \textbf{Otros lenguajes y herramientas:} C, Go, Python, Haskell, Prolog; Linux, Bash, PostgreSQL
\end{itemize}

\section*{Educación e Idiomas}
\textbf{Licenciatura en Ciencias de la Computación}~--- Universidad de Buenos Aires, en curso (53\% completado)

\noindent \textbf{Bachillerato Bilingüe en Ciencias e Informática}~--- Liceo Franco Argentino Jean Mermoz

\noindent \textbf{Idiomas:} Español: nativo \;\textperiodcentered\; Inglés: avanzado \;\textperiodcentered\; Francés: avanzado
